% =========================================================================
% English translation (generated by AI using Claude Opus 4.8 (Max effort)).
% This file was translated from Hungarian to English by AI.
% DISCLAIMER: Please review against the original Hungarian .tex files, before relying on it!
% SOURCE: 4. Számelmélet, gráfok.tex
% =========================================================================
\documentclass[margin=0px]{article}

\usepackage{listings}
\usepackage[utf8]{inputenc}
\usepackage{graphicx}
\usepackage{float}
\usepackage[a4paper, margin=0.7in]{geometry}
\usepackage{subcaption}
\usepackage{amsthm}
\usepackage{amssymb}
\usepackage{amsmath}
\usepackage{fancyhdr}
\usepackage{setspace}

\onehalfspacing

\renewcommand{\figurename}{Figure}
\newenvironment{tetel}[1]{\paragraph{#1 \\}}{}

\newcommand{\N}{\mathbb{N}}
\newcommand{\Z}{\mathbb{Z}}
\newcommand{\R}{\mathbb{R}}
\newcommand{\Q}{\mathbb{Q}}
\newcommand{\C}{\mathbb{C}}

\pagestyle{fancy}
\lhead{\it{CS BSc Final Exam Topics}}
\rhead{4. Number theory, graphs}

\title{\textbf{{\Large ELTE IK - Computer Science BSc} \vspace{0.2cm} \\ {\huge Final Exam Topics}} \vspace{0.3cm} \\ 4. Number theory, graphs}
\author{}
\date{}

\begin{document}
\maketitle

\begin{center}
\fbox{\parbox{0.92\textwidth}{\small \textit{Note: This document was translated from Hungarian to English by AI. Please verify the information against the original Hungarian source before relying on it.}}}
\end{center}

\begin{tetel}{Number theory, graphs, coding theory}
    Sets, relations, functions and operations. Complex numbers. Enumerations on finite sets. Undirected and directed graphs, trees, Eulerian and Hamiltonian graphs, data structures of graphs. Basic concepts of number theory, divisibility, congruence, primes. Polynomials and their operations, division with remainder.
\end{tetel}
\section{Number theory}
\subsection{Sets}
A set (system, class, collection, ...) is the conceptual envelope of its elements. A set is uniquely determined by its elements.
\begin{description}
    \item[Basic concepts] \hfill
        \begin{itemize}
            \item Empty set \\
                  The set that has no elements is the empty set. Notation: $\emptyset$. Based on the axiom of determinacy this is unique
            \item Subset \\
                  We say that \textbf{A is a subset of B} ($A \subseteq B$) if $\forall a \in A : a \in B$, that is, B contains every element of A.
                  \textbf{A is a proper subset of B} \(A \subset B\) if $A \subseteq B$, but $A \neq B$, that is, B has at least one element that is not an element of A.
            \item Power set \\
                  If {\it A} is a set, then the set system whose elements are exactly the subsets of the set {\it A} is called the power set of {\it A}, and is denoted by $2^A$.
                  \begin{itemize}
                      \item $A = \emptyset , 2^\emptyset = {\emptyset}$
                      \item $A = \{a\} , 2^{\{a\}} = \{\emptyset, \{a\}\}$
                      \item $A = \{a, b\} , 2^{\{a,b\}} = \{\emptyset, \{a\}, \{b\}, \{a, b\}\}$
                      \item $|2^A| = 2^{|A|}$
                  \end{itemize}
        \end{itemize}
    \item[Operations] \hfill
        \begin{itemize}
            \item Union \\
                  The union of the sets {\it A} and {\it B}: $A \cup B$ is the set that contains exactly the elements of {\it A} and {\it B}.
            \item Intersection \\
                  The intersection of the sets {\it A} and {\it B}: $A \cap B$ is the set that contains exactly the common elements of {\it A} and {\it B}: $A \cap B = \{ x \in A : x \in B \}$
            \item Disjoint \\
                  If $A \cap B = \emptyset$, then {\it A} and {\it B} are disjoint.
            \item Difference \\
                  The difference of the sets {\it A} and {\it B} is $A \diagdown B = \{ x \in A : x \notin B \}$
            \item Complement \\
                  For a fixed universe {\it X} and a subset $A \subseteq X$, the complement of the set {\it A} is $ \overline{A} = A' = X \diagdown A$
            \item Symmetric difference \\
                  $A \triangle B = (A \diagdown B) \cup (B \diagdown A)$
        \end{itemize}
    \item[Properties] \hfill
        \begin{itemize}
            \item Union
                  \begin{itemize}
                      \item $A \cup \emptyset = A$
                      \item $A \cup ( B \cup C ) = (A \cup B) \cup C$ (associativity)
                      \item $A \cup B = B \cup A$ (commutativity)
                      \item $A \cup A = A$ (idempotence)
                      \item $A \subseteq B \Longleftrightarrow A \cup B = B$
                  \end{itemize}
            \item Intersection
                  \begin{itemize}
                      \item $A \cap \emptyset = \emptyset$
                      \item $A \cap ( B \cap C ) = (A \cap B) \cap C$ (associativity)
                      \item $A \cap B = B \cap A$ (commutativity)
                      \item $A \cap A = A$ (idempotence)
                      \item $A \subseteq B \Longleftrightarrow A \cap B = A$
                  \end{itemize}
            \item Distributivity properties of union and intersection
                  \begin{itemize}
                      \item $A \cap ( B \cup C ) = (A \cap B) \cup (A \cap C)$
                      \item $A \cup ( B \cap C ) = (A \cup B) \cap (A \cup C)$
                  \end{itemize}
            \item Difference
                  \begin{itemize}
                      \item $A \diagdown B = A \cap \overline{B}$
                  \end{itemize}
            \item Complement
                  \begin{itemize}
                      \item $\overline{\overline{B}} = A$
                      \item $\overline{\emptyset} = X$
                      \item $\overline{X} = \emptyset$
                      \item $A \cap \overline{A} = \emptyset$
                      \item $A \cup \overline{A} = X$
                      \item $A \subseteq B \Longleftrightarrow \overline{B} \subseteq \overline{A}$
                      \item $\overline{A \cap B} = \overline{A} \cup \overline{B}$
                      \item $\overline{A \cup B} = \overline{A} \cap \overline{B}$
                  \end{itemize}
            \item Symmetric difference
                  \begin{itemize}
                      \item $A \triangle B = (A \cup B) \diagdown (B \cap A)$
                  \end{itemize}
        \end{itemize}
\end{description}

\subsection{Relations, orderings}
\begin{description}
    \item[Basic concepts] \hfill
        \begin{itemize}
            \item Ordered pair \\
                  $(x,y)$ is an ordered pair if $(x,y) = (u,v) \Longleftrightarrow x = u \ \land \ y = v$. We define this property with sets:
                  \[ (x,y) := \{ \{x\}, \{x, y\} \} \]
            \item Cartesian product \\
                  The Cartesian product or direct product of the sets $X,Y$:
                  \[ X \times Y := \{ (x,y) : x \in X, y \in Y \} \]
            \item Binary relation \\
                  A set is called a binary relation if every element of it is an ordered pair.
                  If $R$ is a binary relation and $(x,y) \in R$, then we often write: $xRy$
            \item Relation \\
                  If for sets $X,Y$ we have $R \subset X\times Y$, then $R$ is a relation between $X$ and $Y$.
            \item Domain \\
                  The domain of the binary relation $R$:
                  \[ \textrm{dmn}(R) := \{ x\  | \ \exists y : (x,y)\in R  \}\]
            \item Range \\
                  The range of the binary relation $R$:
                  \[ \textrm{rng}(R) := \{ y\  | \ \exists x : (x,y)\in R  \}\]
            \item Inverse \\
                  The inverse of a binary relation $R$:
                  \[R^{-1} := \{(a,b) : (b,a) \in R\} \]
            \item Image of a set \\
                  Let $R$ be a binary relation, and $A$ a set. The image of the set $A$:
                  \[R(A) := \{y \ | \ \exists x\in A: (x,y) \in R\} \]
            \item Composition \\
                  The composition of the binary relations $R$ and $S$:
                  \[ R \circ S := \{ (x,y) \ | \ \exists z : (x,z) \in S \ \land \ (z,y) \in R \ \} \]
        \end{itemize}
    \item[Properties] \hfill \\
        $R$ is a binary relation on $X$ (that is, $R \subset X\times X$)
        \begin{enumerate}
            \item transitive \\
                  \[\forall x,y,z : (x,y)\in R \ \land \ (y,z) \in R \Longrightarrow (x,z) \in R \]
            \item symmetric \\
                  \[\forall x,y : (x,y)\in R \Longrightarrow (y,x) \in R \]
            \item antisymmetric \\
                  \[\forall x,y : (x,y)\in R \ \land \ (y,x) \in R \Longrightarrow x = y\]
            \item strictly antisymmetric \\
                  \[\forall x,y : (x,y)\in R \Longrightarrow (y,x) \notin R\]
            \item reflexive \\
                  \[\forall x \in X : (x,x)\in R\]
            \item irreflexive \\
                  \[\forall x \in X : (x,x)\notin R\]
            \item trichotomous \\
                  If for every $x,y \in X$ exactly one of the following holds
                  \begin{enumerate}
                      \item[a)] $x=y$
                      \item[b)] $(x,y) \in R$
                      \item[c)] $(y,x) \in R$
                  \end{enumerate}
            \item dichotomous \\
                  \[\forall x,y \in X : (x,y) \in R \ \lor \ (y,x) \in R\]
                  In other words, the elements are comparable.
        \end{enumerate}
    \item[Orderings] \hfill
        \begin{itemize}
            \item Equivalence relation, classification \\
                  $X$ a set, $R$ a binary relation on $X$, is an equivalence relation if it is
                  \begin{itemize}
                      \item Reflexive
                      \item Transitive
                      \item Symmetric
                  \end{itemize}

                  A system $\mathcal{O}$ of subsets of $X$ is called a classification (partition) if $\mathcal{O}$ is a set system consisting of pairwise disjoint nonempty sets for which $\cup\mathcal{O} = X$
                  \\\\
                  Theorem: \\
                  An equivalence relation determines a classification. Conversely: for a classification $\mathcal{O}$, \\ ${R = \cup\{Y\times Y : Y \in \mathcal{O} \}}$ is an equivalence relation.

            \item Partial order\\
                  $X$ a set, $R$ a binary relation on $X$, is a partial order if it is
                  \begin{itemize}
                      \item Reflexive
                      \item Transitive
                      \item Antisymmetric
                  \end{itemize}
            \item Total order \\
                  $X$ a set, $R$ a binary relation on $X$, is a (total) order if it is
                  \begin{itemize}
                      \item Reflexive
                      \item Transitive
                      \item Antisymmetric
                      \item Dichotomous
                  \end{itemize}
                  In plain words, if a partial order is dichotomous (that is, all of its elements are comparable), then it is a (total) order.
            \item Strict and weak relation, ordering \\
                  $X$ a set, $R$,$S$ relations on $X$. If
                  \[ xRy \ \land \ x \neq y \Rightarrow xSy\]
                  then we call $S$ the strictification of $R$.\\
                  Conversely, if
                  \[  xRy \ \lor \ x = y \Rightarrow xTy \]
                  then $T$ is the weak relation corresponding to $R$.\\

                  \textit{Remark: This is essentially about removing and adding reflexivity. For a partial order the corresponding strict relation (strict partial order) is therefore irreflexive, and consequently also strictly antisymmetric. Conversely: The weak relation corresponding to a strict partial order on X (trans., irrefl., strictly antisym.) is a partial order. }
        \end{itemize}
    \item[Bounds] \hfill
        \begin{itemize}
            \item Least, greatest, minimal, maximal element \\
                  By the least (first) element of a partial order ($\preccurlyeq$) on a set $X$ we mean an element $x\in X$ for which: $\forall y \in X : x\preccurlyeq y$. (Such an element need not exist, but if it does, it is unique).\\
                  Similarly, the greatest (last) element is an $x\in X$ such that $\forall y \in X : y \preccurlyeq x$.\\

                  We call $x$ minimal if there is no element smaller than it, and maximal if there is no element greater than it. (In contrast to least/greatest elements, there can be several minimal/maximal elements. If, however, $X$ is ordered, then least=minimal, greatest=maximal.)
            \item Lower bound, upper bound \\
                  $X$ a partially ordered set, $Y \subset X$. The element $x \in X$ is a lower bound of $Y$ if $\forall y \in Y : x \preccurlyeq y$. (upper bound: $\forall y \in Y : y \preccurlyeq x$). It can be seen that $x$ is not necessarily an element of $Y$; moreover it is also possible that $Y$ has no lower/upper bound, or even several. If, however, $x\in Y$, then it is unique and it is the least element of Y.
            \item Infimum, supremum \\
                  If among the lower bounds there is a greatest element, it is called the infimum of $Y$. (Notation: inf$Y$) \\
                  If among the upper bounds there is a least element, it is called the supremum of $Y$. (Notation: sup$Y$)
            \item Lower bound, upper bound property \\
                  $X$ a partially ordered set. If $ \forall \emptyset \neq Y \subset X$ : $Y$ is bounded above and has a supremum, then it has the upper bound property.
                  Likewise, if $ \forall \emptyset \neq Y \subset X$ : $Y$ is bounded below and has an infimum, then it has the lower bound property.
        \end{itemize}
\end{description}
\subsection{Functions and operations}
\subsubsection{Functions}
\begin{description}
    \item[Definition] \hfill \\
        A relation $f$ is a function if
        \[ (x,y) \in f \ \land (x,y') \ \in f \Longrightarrow y = y' \]
        In other words, for every $x$ there exists at most one $y$ such that $(x,y) \in f$

        Thus for every $x \in$ dmn($f$) we have $f(x) = \{y\}$, which we also usually denote $f(x) = y$ or $f: x \mapsto y$ or $f_x = y$.

    \item[Domain, range] \hfill \\
        We use the notation $f : X \rightarrow Y$ if dmn($f$) $ = X $. \\
        We use the notation $f \in X \rightarrow Y$ if dmn($f$) $\subset X$ (when dmn($f$)$ \subsetneq X$ can also occur).\\
        In both cases rng($f$) $\subset Y$.

    \item[Injective] \hfill \\
        a function $f$ is one-to-one/injective if
        \[ f(x) = y \ \land \ f(x') = y \ \Longrightarrow \ x = x' \]
        This is equivalent to $f^{-1}$ also being a function.
    \item[Surjective] \hfill \\
        The function $f$ is surjective if
        \[ \forall y \in Y : \exists x\in X : f(x) = y\]
        That is, rng($f$) = $Y$. In plain words the function $f$ maps onto the whole of $Y$.
    \item[Bijective] \hfill \\
        If the function $f$ is injective and surjective, then it is bijective.
    \item[Indexed family] \hfill \\
        The value of the function $x$ taken at position $i$ is also usually denoted $x_i$. In this case we often call the domain dmn($f$) = $I$ the index set, its elements indices, rng($f$) the indexed set, and the function $x$ itself an indexed family.
\end{description}
\subsubsection{Operations}
\begin{description}
    \item[Definitions] \hfill
        \begin{itemize}
            \item Binary operation \\
                  On a set $X$ a function $f : X \times X \rightarrow X$ is a binary operation.
            \item Unary operation \\
                  On a set $X$ a function $f : X \rightarrow X$ is a unary operation.
            \item Nullary operation \\
                  $X$ a set, $f : \{\emptyset \} \rightarrow X $ a nullary operation. (Essentially the selection of an element)
        \end{itemize}
    \item[Properties] \hfill
        \begin{itemize}
            \item Let $\spadesuit, \copyright$ be binary operations on $X$.
                  \begin{enumerate}
                      \item $\spadesuit$ is associative if
                            \[ \forall x,y,z \in X : \quad (x \ \spadesuit \ y ) \ \spadesuit \ z = x \  \spadesuit \  (y \ \spadesuit z) \]

                      \item $\spadesuit$ is commutative if
                            \[ \forall x,y \in X : \quad x \ \spadesuit \ y = y \  \spadesuit \  x \]

                      \item $\spadesuit$ is distributive over $\copyright$ if $\forall x,y,z \in X$:
                            \[ x \  \spadesuit \  (y \ \copyright \ z) = (x \  \spadesuit \ y) \ \copyright \ (x \ \spadesuit \ z) \quad \textrm{ - left} \]
                            \[ (y \ \copyright \ z) \  \spadesuit \ x = (y \  \spadesuit \ x) \ \copyright \ (z \ \spadesuit \ x) \quad \textrm{ - right} \]
                  \end{enumerate}

            \item Let $\heartsuit$ be a binary operation on $X$ and $\S$ a binary operation on $Y$.
                  $f : X \rightarrow Y$ is operation-preserving if:
                  \[ \forall x_1,x_2 \in X:  f(x_1 \heartsuit x_2) = f(x_1) \ \S \ f(x_2) \]
        \end{itemize}
\end{description}
\subsection{Number concept, complex numbers}
\subsubsection{Number concept}
\begin{description}
    \item[Algebraic structures] \hfill
        \begin{enumerate}
            \item Groupoid \\
                  A set $G$ with an operation $\star$, that is, the pair $(G, \star)$, is called a groupoid.
            \item Semigroup\\
                  If in a groupoid the operation $\star$ is associative, then the groupoid is a semigroup.
            \item Monoid \\
                  A semigroup with an identity element is called a monoid.\\
                  \textit{Remark: $ a \in G$ is an identity element if $\forall g \in G : a \star g = g \star a = g$}
            \item Group \\
                  If in a monoid every element has an inverse, then we speak of a group. \\
                  \textit{Remark: $g,g^{-1} \in G$ and $\xi\in G$ is the identity element, then $g^{-1}$ is the inverse of $g$ if $g\star g^{-1} = \xi$ and $g^{-1} \star g = \xi$}
            \item Abelian group \\
                  If in a group the operation is commutative, then it is an Abelian group.
            \item Ring \\
                  $(R,+,\cdot )$ is a ring if it is an Abelian group under addition, a semigroup under multiplication, and both-sided distributivity holds.

                  If multiplication is commutative, then it is a commutative ring.

                  If multiplication has an identity element, then it is a ring with identity.
            \item Integral domain \\
                  A commutative ring without zero divisors.

                  \textit{Zero divisor: $x,y$ are distinct nonzero elements, but $x\cdot y = 0$}
            \item Ordered integral domain \\
                  An integral domain $R$ is an ordered integral domain if it is an ordered set, and moreover addition and multiplication are monotone.

                  \textit{Addition monotone: $x,y,z \in R$ and $x \leq y \ \Rightarrow \ x+z \leq y+z$ \\
                      Multiplication monotone: $x,y \in R$ and $x,y\geq0 \ \Rightarrow \ x\cdot y \geq 0$ }
            \item Field \\
                  A ring $R$ is a field if $R\setminus\{0\}$ is an Abelian group under multiplication.
            \item Ordered field \\
                  If a field is an ordered integral domain, then it is an ordered field.
        \end{enumerate}
        \item[Natural numbers]\hfill
        \begin{itemize}
            \item Peano axioms \\
                  Let $\N$ be a set and $^+$ a function defined on $\N$. The following conditions are called the Peano axioms:
                  \begin{enumerate}
                      \item $0 \in \N \qquad$ - $0$ is a nullary operation on $\N$
                      \item if $n \in \N$, then $n^+ \in N \qquad$ - $^+$ is a unary operation on $\N$
                      \item if $n \in \N$, then $n^+ \neq 0 \qquad $ - $0$ is not in the range of $^+$
                      \item if $n,m \in \N$, and $m^+ = n^+$, then $n = m \qquad $ - $^+$ is injective
                      \item if $S \subset \N, 0 \in S$, and $ n \in S : n^+\in S$, then $S = \N \qquad $ - the principle of mathematical induction
                  \end{enumerate}
            \item Operations
                  \begin{itemize}
                      \item addition \\
                            $k,m,n \in \N$, then:
                            \begin{enumerate}
                                \item $(k+m)+n = k+(m+n)$ \textit{- associativity}
                                \item $n+0 = 0+n = n$ \textit{- 0 is the zero element (additive identity)}
                                \item $n+k = k+n$ \textit{- commutativity}
                                \item $n+k = m+k$ or $k+n = k+m$, then $m=n$ \textit{- cancellation rule}
                            \end{enumerate}
                      \item multiplication \\
                            $k,m,n \in \N$, then:
                            \begin{enumerate}
                                \item $(k\cdot m)\cdot n = k\cdot (m\cdot n)$ \textit{- associativity}
                                \item $ 0\cdot n = n\cdot 0 = 0$
                                \item $n\cdot 1 = 1\cdot n = n$ \textit{- 1 is the identity element (multiplicative identity)}
                                \item $n\cdot k = k\cdot n$ \textit{- commutativity}
                                \item $k\cdot (m+n) = k\cdot m + \cdot n$, and $(m+n) \cdot k = m\cdot k+n\cdot k$ \textit{- distributivity}
                                \item for $k\neq 0$: $n\cdot k = m\cdot k$, then $m=n$ \textit{- cancellation rule}
                            \end{enumerate}
                  \end{itemize}
        \end{itemize}
    \item[Integers] \hfill \\
        In the realm of natural numbers, with respect to addition only zero has an inverse; in other words, subtraction generally cannot be performed.

        Consider the relation $\sim\ \subset\N\times\N$ for which $(m,n) \sim (m',n')$ if $m+n' = m'+n$. And take the addition $(m,n)+(m',n') = (m+m',n+n')$. The relation $\sim$ is an equivalence relation; we denote the set of equivalence classes by $\Z$. We call the elements of $\Z$ integers.

        Addition is compatible with the equivalence, so it is defined among the integers, and $(\Z, +)$ is an Abelian group.

        So $(\Z, +, \cdot)$ is a ring.

        \textit{Remark: an operation $*$ is compatible with the equivalence relation $\asymp$ if the following holds: $ x \asymp x' \ \land \ y \asymp y' \ \Longrightarrow \ x * y \asymp x'*y'$}
    \item[Rational numbers] \hfill \\
        In the realm of integers, among the nonzero elements only $1$ and $-1$ have a multiplicative inverse; in other words, division generally cannot be performed.

        Consider on $\Z \times (\Z\setminus\{0\})$ the relation $\sim$ for which $(m,n) \sim (m',n')$ if $mn' = nm'$. And take the addition $(m,n)+(m',n') = (mn'+nm', nn')$ and the multiplication $(m,n)\cdot(m',n')=(mm', nn')$. The relation $\sim$ is an equivalence relation; we denote the set of equivalence classes by $\Q$. We call the elements of $\Q$ rational numbers.

        $(\Q, +, \cdot)$ is an ordered field.
    \item[Real numbers] \hfill \\
        There is no number $a \in \Q$ whose square is 2. So not every number can be written in the form m/n ($m,n \in \N^+$).

        Archimedean ordering:\\
        An ordered field $F$ is Archimedean ordered if $x,y\in F: \exists n \in \N : nx \geq y \quad (x>0)$

        The ordered field of rational numbers is Archimedean ordered, but does not have the upper bound property.

        An ordered field with the upper bound property is called the field of real numbers, and is denoted $\R$. ($\exists!\R$)
\end{description}
\subsubsection{Complex numbers}
The need for complex numbers arose from the Cardano formula for solving cubic equations. Namely, in the case when the equation has three distinct real roots, a negative number appears under the root sign in the formula. The rules for computing with ``imaginary'' numbers, and the connection with trigonometric functions, were gradually clarified.

\begin{description}
    \item[Definition] \hfill \\
        The set of complex numbers is $\C = \R \times \R$. $\C$ is a field with the addition $(x,y)+(x',y') = (x+x', y+y')$ and the multiplication $(x,y)\cdot(x',y') = (xx'-yy', y'x+yx')$. The set of complex numbers is not an ordered field, since (by a theorem) in an ordered integral domain $ x \neq 0 \ \Rightarrow x^2 > 0$. (This, however, does not hold for $(0,1)^2=i^2 = -1$).

            [In the realm of complex numbers $(0,0)$ is the zero element, $(1,0)$ the identity element, the additive inverse of $(x,y)$ is $(-x,-y)$, and the multiplicative inverse of the pair $(0,0) \neq (x,y)$ is the pair $(\frac{x}{x^2+y^2}, \frac{-y}{x^2+y^2})$.]
    \item[Identification of real numbers] \hfill \\
        Since $(x,0)+(x',0) = (x+x',0)$ and $(x,0)\cdot(x',0) = (xx',0)$, we can identify all complex numbers $(x,0), x\in\R$ with $\R$.
    \item[Algebraic form of complex numbers] \hfill \\
        Since
        \[(x,y) = (x,0)+(y,0)\cdot i = x+yi\]
        we can also write complex numbers in the algebraic form $a+bi$.

        Then the real number Re($z$) = $x$ is called the real part of the complex number $z = (x,y)$, and the real number Im($z$) = $y$ its imaginary part.
    \item[Conjugate] \hfill \\
        The conjugate of the complex number $ z = x+yi$: $\overline{z} = x-yi$

        Its properties:
        \begin{enumerate}
            \item $\overline{z+w} = \overline{z}+\overline{w}$
            \item $\overline{z\cdot w} = \overline{z}\cdot\overline{w}$
            \item $\overline{\overline{z}} = z$
            \item $z + \overline{z}$ = 2Re($z$)
            \item $z - \overline{z}$ = $i\cdot2$Im($z$)
        \end{enumerate}
    \item[Absolute value] \hfill \\
        The absolute value of the complex number $z=(x,y)$: $|z| = \sqrt{x^2+y^2}$

        Its properties:
        \begin{enumerate}
            \item $z\cdot\overline{z} = {|z|}^2$
            \item $\frac{1}{z}= \frac{\overline{z}}{|z|^2} $
            \item $|z| = \overline{|z|}$
            \item $|z\cdot w | = |z|\cdot|w|$
            \item $|z+w| \leq |z| +|w|$
        \end{enumerate}
    \item[Trigonometric form] \hfill
        \begin{itemize}
            \item Argument \\
                  For $z \neq 0$ the argument of $z$ is any $t \in\R$ for which Re($z$) = $|z|$cos($t$), and Im($z$) = $|z|$sin($t$). In other words the argument of $z$ is the angle between the vector pointing from the origin to $z$ and the positive real axis.
            \item Trigonometric form \\
                  The trigonometric form of the complex number $z$: $ z = |z|($cos($t$)+$i\cdot$sin($t$)
            \item De Moivre identities \\
                  Let $z = |z|($cos($t$)+$i\cdot$sin($t$)), and $w = |w|($cos($s$)+$i\cdot$sin($s$)). Then
                  \[z\cdot w = |z||w|(\textrm{cos}(t+s)+i\cdot \textrm{sin}(t+s))\]
                  \[\frac{z}{w} = \frac{|z|}{|w|}(\textrm{cos}(t-s)+i\cdot \textrm{sin}(t-s)) \quad (w \neq 0)\]
                  \[ z^n = |z|^n(\textrm{cos}(nt)+i\cdot \textrm{sin}(nt)) \quad (n \in \Z)\]
            \item Root extraction \\
                  Let $z^n = w$, then:
                  \[ \sqrt[n]{w} = \Bigg\{z_k = \sqrt[n]{|w|}\bigg(\textrm{cos}\Big(\frac{t+2k\pi}{n}\Big)+\textrm{sin}\Big(\frac{t+2k\pi}{n}\Big)\bigg), k=0,...,n-1\Bigg\} \]
                  But since this notation can be confused with the (made unambiguous) real root extraction among reals, we do not use this notation. Instead, let us introduce the concept of the $n$-th complex root of unity:
                  \[\varepsilon_k = \textrm{cos}\bigg(\frac{2k\pi}{n}\bigg)+i\cdot\textrm{sin}\bigg(\frac{2k\pi}{n}\bigg), \quad k=0,...,n-1\]
                  After this we can obtain the roots of $w$ with the help of $z$ and the $n$-th complex roots of unity: $z\varepsilon_0, ..., z\varepsilon_{n-1}$
        \end{itemize}

\end{description}
\subsection{Enumerations on finite sets}
\begin{description}
    \item[Finite sets] \hfill
        \begin{itemize}
            \item Equivalence of sets \\
                  Sets $X,Y$ are equivalent if there exists a bijection mapping $X$ onto $Y$.\\
                  Notation: $X \sim Y$
            \item Finite and infinite sets \\
                  A set $X$ is finite if $\exists n\in\N : X \sim \{1,2,...,n\}$, otherwise infinite. If such an $n$ exists, then it is unique, and then it is called the cardinality of the set. Notation: $ \#(X)$
        \end{itemize}
    \item[Pigeonhole principle] \hfill \\
        If $X,Y$ are finite sets and $\#(X) > \#(Y)$, then a map $f:X\rightarrow Y$ cannot be one-to-one (that is, a bijection).
    \item[Enumerations] \hfill
        \begin{itemize}
            \item Permutation \\
                  A permutation of a set $A$ is a one-to-one map of it onto itself. The number of all permutations of the set $A$:
                  \[P_n = \prod\limits_{k=1}^{n} k \ = \ n!\]
            \item Variation \\
                  The sequences $a_1,a_2,...,a_k$ of distinct elements that can be formed from the elements of the set $A$ are called the variations of class $k$ of the set $A$. If $A$ is finite ($\#(A) = n$), then the number $V_n^k$ equals the number of one-to-one maps from $\{1,2,...,k\}$ into $\{1,2,...,n\}$:
                  \[ V_n^k = \frac{n!}{(n-k)!}\]
            \item Combination \\
                  The subsets of $k\in\N$ elements of a set $A$ are called the combinations of class $k$. If $A$ is finite, then the number $C_n^k$ equals the number of $k$-element subsets of $\{1,2,...,n\}$.
                  \[C_n^k = \dbinom{n}{k} = \frac{n!}{k!(n-k)!}\]
            \item Permutation with repetition \\
                  The repetitions of the elements of the set $A = \{a_1,\dotsc,a_r\}$ are $i_1,\dotsc,i_r$. (The permutations with repetition of the elements are sequences of $i_1+\cdots+i_r = n$ terms, in which the element $a_j$ occurs $i_j$ times.)
                  \[P_n^{i_1,\dotsc,i_r} = \frac{n!}{i_1!i_2!\cdots i_r!}\]
            \item Variation with repetition \\
                  The (not necessarily distinct) sequences $a_1,\cdots,a_k$ that can be formed from the elements of the finite set $A$ are called the variations with repetition of class $k$ of the set $A$.
                  \[^iV_n^k = n^k\]
            \item Combination with repetition \\
                  Let $A$ be a finite set. Selecting $k$ elements from the set, allowing repetitions but disregarding order, we obtain the combinations with repetition of class $k$ of the set $A$.
                  \[^iC_n^k = \dbinom{n+k-1}{k} \]
        \end{itemize}
    \item[Theorems] \hfill \\
        \begin{itemize}
            \item Binomial theorem \\
                  $x,y \in R$ (commutative ring with identity), $n\in\R$. Then
                  \[(x+y)^n = \sum\limits_{k=0}^n\dbinom{n}{k}x^ky^{n-k} \]
            \item Multinomial theorem \\
                  $r,n\in\N$ and $x_1, x_2, \cdots, x_r \in R$ (commutative ring with identity), then
                  \[(x_1+\cdots+x_r)^n = \sum\limits_{i_1+\cdots+i_r = n}P_n^{i_1,\cdots,i_r}x_1^{i_1}x_2^{i_2}\cdots x_r^{i_r} \qquad (i_1,\cdots,i_r \in\N)\]
            \item Inclusion-exclusion (sieve) formula \\
                  $X_1,\cdots,X_k \subset X$ (finite set). $f$ a function defined on $X$, mapping into an Abelian group. Let:
                  \[S=\sum\limits_{x\in X}f(x)\]
                  \[ S_r = \sum\limits_{1\leq i_1 \leq \cdots \leq i_r \leq k}\Bigg(\sum\limits_{x \in X_{i_1} \cap \cdots \cap X_{i_r}}f(x)\Bigg) \]
                  and
                  \[S_0 = \sum\limits_{x\in X \setminus \cup_{i=1}^k Xi}f(x)\]
                  Then
                  \[S_0 = S - S_1+S_2-S_3+\cdots+(-1)^kS_k \]
        \end{itemize}
\end{description}
\subsection{Basic concepts of number theory, division with remainder, linear congruence equations}
\subsubsection{Basic concepts of number theory}
\begin{description}
    \item[Divisibility in an integral domain with identity] \hfill \\
        $R$ an integral domain with identity, $a,b\in R$. If $\exists c\in R: a = bc$, then $b$ is a divisor of $a$ ($a$ is a multiple of $b$). Notation: $b|a$

        Except for $b = 0$, at most one such $c$ exists.

        Properties of divisibility in an integral domain with identity.
        \begin{itemize}
            \item If $b|a$ and $b'|a'$, then $bb'|aa'$
            \item $\forall a\in R: a|0$ (every element divides zero)
            \item $0|a \Leftrightarrow a = 0$ (zero divides only itself)
            \item $\forall a\in R: 1|a$ (the identity element divides every element)
            \item $b|a \Rightarrow \forall c \in R : bc|ac$
            \item $bc|ac$ and $c\neq0 \Rightarrow b|a$
            \item $b|a_i$ and $c_i \in R,\ (i = 1,\cdots,j) \ \Rightarrow \ b|\sum_{i=1}^{j}a_ic_i$
            \item the $|$ relation is reflexive and transitive
        \end{itemize}
    \item[Irreducible element and prime element] \hfill \\
        $0,1 \neq a \in R$ is irreducible if, for $a = bc$, $b$ or $c$ is a unit ($b,c \in R$).

        $0,1 \neq p \in R$ is prime if $\forall a,b \in R : p|ab$ implies $p|a$ or $p|b$
    \item[Greatest common divisor, least common multiple, relatively prime] \hfill \\
        $R$ an integral domain with identity. $b\in R$ is a greatest common divisor of the elements $a_1,\cdots,a_n \in R$ if $b|a_i$ and ($b'|a_i$ implies $b'|b$). If $b$ is a unit, then $a_1, \cdots , a_n$ are relatively prime.

        $b\in R$ is a least common multiple of the elements $a_1,\cdots,a_n \in R$ if $a_i|b$ and ($a_i|b'$ implies $b|b'$).
    \item[Extended Euclidean algorithm] \hfill \\
        The procedure determines the greatest common divisor ($d\in\Z$) of the numbers $a,b \in \Z$, as well as numbers $x,y \in\Z$ such that $d = ax+by$
    \item[Fundamental theorem of arithmetic] \hfill \\
        Every positive natural number can be uniquely (up to order) decomposed as a product of prime numbers.
    \item[Sieve of Eratosthenes] \hfill \\
        To determine the primes up to a given $n$:
        Write down the numbers from 2 to $n$. The first number (2) is prime, all its multiples are composite, cross these out. Among the remaining numbers, the first (3) is likewise prime, etc. At the end of the procedure the primes not greater than $n$ remain.
\end{description}
\subsubsection{Division with remainder}
Let $R$ be an integral domain with identity, $f,g\in R[x], g \neq 0$ and suppose that the leading coefficient of $g$ is a unit in $R$. Then
\[ \exists! q,r \in R[x] : f = g\cdot q+r \qquad (\text{deg}(r) < \text{deg}(g) )\]

\subsubsection{Horner's scheme}
Horner's method is suitable for computing the value of a polynomial at a substitution point. (Together with this, of course, one can also decide whether a given value $c$ is a root of the polynomial or not. Above degree 4 there is not even an analytic solution for this.)

The essence of the method is that computing the value of the polynomial $f_nx^n+f_{n-1}x^{n-1}+\cdots+f_0$, which would otherwise require a great many multiplications and additions. By transforming the polynomial, however, we can reduce the number of operations. Applying division with remainder:
\[f_nx^n+f_{n-1}x^{n-1}+\cdots+f_0 = (f_nx^{n-1}+f_{n-1}x^{n-2}+\cdots)x+f_0\]
Continuing this recursively we arrive at the following form:
\[(((f_nx+f_n-1)x+f_n-2)x+\cdots)x+f_0\]
The computation of the substitution value can be carried out more easily in a table.

\begin{center}
    \begin{tabular}{|c|c|c|c|c|c|}
        \hline     & $f_n$ & $f_{n-1}$      & $f_{n-2}$                 & $\cdots$ & $f_0$  \\
        \hline $c$ & $f_n$ & $f_nc+f_{n-1}$ & $(f_nc+f_{n-1})c+f_{n-2}$ & $\cdots$ & $f(c)$ \\
        \hline
    \end{tabular}
\end{center}
Filling in the table proceeds as follows:
\begin{enumerate}
    \item In the first row we write the coefficients of the polynomial
    \item In the first cell of the second row we write the value of the argument.
    \item Under the leading coefficient we write it itself.
    \item We continue by filling in the cells of the second row
    \item \label{itm:horner_recursive} We multiply the element of the previous cell by the argument
    \item To the product we add the current coefficient
    \item We write the sum into the current cell
    \item We continue with step \ref{itm:horner_recursive} until we reach the last cell
\end{enumerate}

\noindent
The substitution value of the polynomial goes into the last cell. (If this is zero, then the argument is a root of the polynomial.)
\subsubsection{Linear congruence equations}
\begin{description}
    \item[Congruence] \hfill \\
        If $a,b,m \in\Z$ and $m|(a-b)$, then we say that $a$ and $b$ are congruent modulo $m$ (Notation: ${a \equiv b \mod{m}}$).

        Congruence is an equivalence relation for any $m$. If $a\in\Z$ then the elements of the equivalence class are of the form $a+km, k\in\Z$.
    \item[Residue classes] \hfill \\
        We call the equivalence classes with respect to the modulus $m\in\Z$ residue classes. We represent residue classes by their elements. (The residue class represented by the element $a$ is $\widetilde{a} \mod{m}$).

        If some element of a residue class is relatively prime to the modulus, then all of them are, and the residue class is called a reduced residue class.

        A system of pairwise incongruent integers is called a residue system.

        If a residue system contains an element from every residue class, then it is a complete residue system.

        If a residue system contains an element exactly from the reduced residue classes, then it is a reduced residue system.
    \item[Euler's $\varphi$ function] \hfill \\
        $m > 0$ an integer. Euler's $\varphi(m)$ function gives the number of reduced residue classes modulo $m$. This obviously equals the number, among the numbers $0,1,\cdots,m-1$, of those relatively prime to $m$.
    \item[Euler--Fermat theorem] \hfill \\
        $m>1$ an integer, $a$ relatively prime to $m$, then:
        \[a^{\varphi(m)}\equiv 1 \mod{m}\]
    \item[Fermat's theorem] \hfill \\
        Let $p$ be prime, and $a\in\Z: p\nmid a$, then
        \[a^{p-1}\equiv 1 \mod p \]
    \item[Solving a linear congruence] \hfill \\
        We look for the solutions of the congruence $ax \equiv b \mod{m}$ ($a,b,m\in\Z$ known). This is equivalent to looking for an $x$ for which (with some $y$) $ax+my = b$.

        Let $d$ = gcd($a,m$). Since $d$ divides $ax+my$, it must also divide $b$, otherwise there is no solution. Thus $\frac{a}{d}x+\frac{m}{d}y = \frac{b}{d}$. Then $a'x+m'y = 1$. With the help of the extended Euclidean algorithm we obtain numbers $u,v$ with which $a'u+m'v = 1$ (since $a', m'$ are relatively prime). Multiplying the equation by $b'$: $a'ub'+m'vb' = b' \Rightarrow x \equiv ub' \mod{m'}$
    \item[Solving a system of linear congruences] \hfill \\
        For two linear congruences the solutions are $x \equiv a \mod{m}$ and $x \equiv b \mod{n}$. For the common solution we have to solve the equation $x = a + my = b+ nz \Leftrightarrow my-nz = b-a$. There is a solution if and only if $d$ = gcd($m,n$) divides $b-a$. Then the solution can be written, with some integer $x_1$, in the form ${ x \equiv x_1 \mod{\textrm{lcm}(m,n)}}$. (For more congruences the procedure can be continued.)
    \item[Chinese remainder theorem] \hfill \\
        $1 < m_1,\cdots,m_n \in\N$ pairwise relatively prime, and $c_1,\cdots,c_n \in\Z$. The system of congruences $x \equiv c_j \mod{m_j}$ ($j=1,\cdots,n$) is solvable, and any two of its solutions are congruent $\mod{m_1m_2\cdots m_n}$
\end{description}
\section{Polynomials and their operations}
\begin{description}
    \item[Definition] \hfill \\
        Let $R$ be a ring. We regard a polynomial as a finite sum of the form $\sum_{i=0}^{n} f_ix^i$, where $n\in\N, f_i \in R$.

        We call the term $f_n$ the leading coefficient of the polynomial.
    \item[Operations] \hfill \\
        Let $R[x]$ be the ring over the infinite sequences $f = (f_0, f_1, \cdots)$ (the ring of polynomials), where $f_i \in R$. Then the operations in $R[x]$:
        \begin{itemize}
            \item Addition: \\
                  \[f+g = (f_0+g_0, f_1+g_1,\cdots)\qquad (f,g \in R[x]) \]
            \item Multiplication: \\
                  \[f\cdot g = h = (h_0,h_1, \cdots) \qquad (f,g,h \in R[x]) \text{, where} \]
                  \[h_k = \sum\limits_{i+j = k}f_ig_j\]
        \end{itemize}

        \textit{Remark: If $R$ is commutative, then so is $R[x]$. If $R$ has identity with the identity element 1, then so does $R[x]$, with the identity element $(1,0,0, \cdots)$.}
\end{description}
\section{Graphs}
\subsection{General and planar graphs}
\begin{description}
    \item[Basic concepts] \hfill
        \begin{itemize}
            \item Undirected graph\\
                  An undirected graph is the ordered triple $G = (V,E, \varphi)$, where:\\
                  $V$ - the set of vertices \\
                  $E$ - the set of edges \\
                  $\varphi$ - incidence relation ($\varphi \in E \times V$)

                  \textit{If $v\in\varphi(e)$, then $v$ is incident to the edge $e$. ($v\in V, e\in E$). An edge always has two ends}

            \item Edge and vertex types
                  \begin{itemize}
                      \item Isolated vertex \\
                            $v\in V$ is an isolated vertex if $\nexists e \in E: v\in \varphi(e)$
                      \item Parallel edge \\
                            edges $e,e'\in E$ are parallel edges if $\varphi(e) = \varphi(e')$
                      \item Loop edge \\
                            $e\in E$ is a loop edge if $|\varphi(e) | = 1$
                  \end{itemize}
            \item Directed graph \\
                  A directed graph is the ordered triple $G = (V,E, \psi)$, where:\\
                  $V$ - the set of vertices \\
                  $E$ - the set of edges \\
                  $\psi$ - incidence relation ($\psi \in E \rightarrow V \times V$)

                  \textit{$\psi(e) = (v,v')$, where $v$ is the start point of the edge $e$, $v'$ its end point.}
        \end{itemize}
    \item[Finite, simple graphs - basic concepts] \hfill
        \begin{itemize}
            \item Simple graph \\
                  a graph $G$ is simple if it does not contain parallel or loop edges
            \item Finite graph
                  a graph $G= (V,E,\varphi)$ is finite if $V,E$ are finite sets.
            \item Adjacency, degree\\
                  Two edges are adjacent if they have a common point.\\
                  Two vertices are adjacent if they have a common edge. \\
                  The number of neighbors of $v\in V$ is the degree of $v$. [Notation: deg($v$) = d($v$)]
            \item $r$-regular graphs\\
                  a graph $G$ is $r$-regular if every vertex has degree $r$
            \item Complete graph\\
                  a graph $G$ is a complete graph if every edge is drawn in, in other words it is $(|V|-1)$-regular. (Notation: $K_{|V|}$)
            \item Bipartite graph\\
                  $G$ is a bipartite graph if $V = V' \cup V''$ and $V'\cap V'' = \emptyset$ (disjoint), and an edge runs only between $V'$ and $V''$.

                  \textit{If, moreover, every edge between $V'$ and $V''$ is drawn in this way, then it is a complete bipartite graph. (Notation: $K_{n,m}$, where $n=|V'|, m=|V''|$)}
            \item Subgraph \\
                  $G = (V,E,\varphi)$ is a subgraph of $G'=(V',G',\varphi')$ if $V\subset V \ \land \ E \subset E' \ \land \ \varphi \subset \varphi'$
            \item Walk, trail, path \\
                  In a graph $G$ a walk of length $n$ from $v$ to $v'$ is a sequence
                  \[v_0,e_1,v_1,\cdots,v_{n-1},e_n,v_n\]
                  for which $v=v_1, v'=v_n$ and $v_{i-1},v_{i} \in \varphi(e_i)$

                  A walk is a trail if every edge appears at most once in the sequence.

                  A trail is a path if every vertex appears at most once in the sequence.

                  A walk/trail/path is closed if its start and end points coincide, otherwise open.
            \item Connected graph
                  A graph is connected if there is a path between any two vertices.

                  \textit{This relation is an equivalence relation, whose equivalence classes are called components.}
            \item Labeled, weighted graph \\
                  $G=(V,E,\varphi, C_e, c_e, C_v,c_v)$ ordered 7-tuple denotes a labeled graph, where $C_e, C_v$ are arbitrary sets, and
                  \[c_e:E\rightarrow C_e\]
                  \[c_v:E\rightarrow C_v\]

                  If $C_e = C_v = \R^+$, then the graph is called a weighted graph, and $w$ is the weight of the vertex/edge. (${w(e)=c_e(e)},\ {w(v) = c_v(v)}$)
        \end{itemize}
    \item[Planarity] \hfill
        \begin{description}
            \item[Concepts] \hfill
                \begin{itemize}
                    \item Planarity \\
                          A graph is planar if it can be drawn so that its edges do not cross each other.
                    \item Topological isomorphism \\
                          Two graphs are topologically isomorphic if by finitely many repetitions of the following step or its reverse we obtain one from the other: We remove a degree-two vertex and connect its neighbors.
                    \item Region \\
                          If a graph $G$ is planar, then the regions are the plane figures bounded by the edges. (The unbounded plane figure is also a region.)
                \end{itemize}
            \item[Theorems] \hfill
                \begin{enumerate}
                    \item Every finite graph can be drawn in $\R^3$.
                    \item A finite graph is planar $\Longleftrightarrow$ it can be drawn on a sphere
                    \item Euler's theorem: \\
                          If the finite graph $G$ is a connected, planar graph, then:
                          \[|E|+2 = |V|+|T|\]
                    \item Kuratowski's theorem: \\
                          A finite graph is planar if and only if it does not contain a subgraph topologically isomorphic to $K_5$ or $K_{3,3}$.
                \end{enumerate}
        \end{description}
\end{description}
\subsection{Trees}
\begin{description}
    \item[Tree]\hfill \\
    A graph is called a tree if it is connected and acyclic.
    \item[Spanning tree] \hfill \\
        $F$ is a subgraph of $G$. If $F$ is a tree and its vertex set equals the vertex set of $G$, then $F$ is called a spanning tree of $G$.
    \item[Theorems] \hfill
        \begin{itemize}
            \item If $G$ is a simple graph, then the following conditions are equivalent:
                  \begin{enumerate}
                      \item $G$ is a tree
                      \item $G$ is connected, but deleting any edge it no longer is
                      \item Between two distinct vertices there is only one path
                      \item $G$ is acyclic, but adding an edge it no longer is
                  \end{enumerate}
            \item If $G$ is a simple finite graph, then the following conditions are equivalent:
                  \begin{enumerate}
                      \item $G$ is a tree
                      \item $G$ has no cycle and has $n-1$ edges
                      \item $G$ is connected and has $n-1$ edges
                  \end{enumerate}
        \end{itemize}
    \item[Directed tree] \hfill \\
        A tree for which: $\exists v \in V : d^-(v) = 0$ and $\forall v'\neq v : d^-(v') = 1$
        (One vertex has in-degree 0, the others 1)

        Further concepts:
        \begin{itemize}
            \item the vertex $r \in V, d^-(r) = 0 $ is called the root
            \item the level of vertex $v'$ is the length of the path $r,v'$
            \item $(v,v')\in\psi(e)$, $v$ is the parent of $v'$, $v'$ the child of $v$.
            \item $v$ is a leaf if $d^+(v)=0$
        \end{itemize}
\end{description}
\subsection{Eulerian and Hamiltonian graphs}
\subsubsection{Eulerian graph}
\begin{description}
    \item[Eulerian trail] \hfill \\
        An Eulerian trail is a trail from $v$ to $v'$ in the graph in which every edge appears. If $v=v'$ then this trail is also usually called an Eulerian circuit. A graph having an Eulerian trail is called an Eulerian graph.
    \item[Theorem] \hfill \\
        In a connected finite graph an Eulerian circuit exists if and only if every vertex has even degree.
\end{description}
\subsubsection{Hamiltonian graph}
A Hamiltonian path is a path from $v$ to $v'$ in the graph that contains every vertex. If $v=v'$ then this path is also usually called a Hamiltonian cycle. A graph having a Hamiltonian path is called a Hamiltonian graph.
\subsection{Data structures of graphs}
For the computer representation of graphs, linked lists or matrices are most often used. Linked lists are more economical for sparse graphs, while matrices are economical for dense graphs.
\begin{description}
    \item[Incidence matrix] \hfill \\
        For a directed graph $G = (V,E, \psi)$ we can represent the graph with the help of a matrix $A = \{0,1,-1\}^{n\times m}$, where $V = \{v_1,\cdots, v_n\}$, and $E = \{e_1,\cdots,e_m\}$. Then the individual elements of the matrix:
        \[ a_{ij} =
            \left\{
            \begin{array}{ll}
                1  & \mbox{if} \ v_i \ \text{is the start point of} \ e_j \\
                -1 & \mbox{if} \ v_i \ \text{is the end point of} \ e_j   \\
                0  & \mbox{otherwise}
            \end{array}
            \right.
        \]

        If $G$ is undirected, then $a_{ij} = |a_{i,j}|$
    \item[Vertex matrix (adjacency matrix)] \hfill \\
        With the above notation, in the directed case $B \in \Z^{n\times n}$, where $b_{ij}$ denotes the number of edges going from $v_i$ to $v_j$.

        If $G$ is undirected, then $b_{ii}$ is the number of loop edges of $v_i$, otherwise $b_{ij}$ is the number of edges between the vertices $v_i$ and $v_j$.
\end{description}

\end{document}
